\section{Môi trường thực nghiệm}

Môi trường phải được ghi nhận ngay trước khi đo để tránh sai khác do cập nhật hệ
thống. Ngoài thông tin trong Bảng~\ref{tab:experiment-devices}, báo cáo cần nêu
model input size, camera/video source, batch size, power mode, clock, tản nhiệt,
thời gian warm-up, số vòng lặp và nhiệt độ môi trường.

\begin{table}[H]
\centering
\caption{Cấu hình thiết bị thực nghiệm}
\label{tab:experiment-devices}
\begin{tabularx}{\textwidth}{L{3.6cm}YY}
\toprule
\textbf{Thành phần} & \textbf{Jetson Nano} & \textbf{Jetson Orin Nano}\\
\midrule
GPU & \cbs{thông tin thực tế} & \cbs{thông tin thực tế}\\
CUDA & \cbs{phiên bản} & \cbs{phiên bản}\\
RAM & \cbs{dung lượng} & \cbs{dung lượng}\\
TensorRT & \cbs{phiên bản} & \cbs{phiên bản}\\
Power Mode & \cbs{chế độ} & \cbs{chế độ}\\
OS/JetPack & \cbs{phiên bản} & \cbs{phiên bản}\\
Storage & \cbs{loại thiết bị} & \cbs{loại thiết bị}\\
Cooling & \cbs{cấu hình} & \cbs{cấu hình}\\
\bottomrule
\end{tabularx}
\end{table}

Mỗi phép đo hiệu năng chạy ít nhất \cbs{số lần lặp} lần sau \cbs{số vòng}
warm-up. Báo cáo sử dụng median và các phần vị p90/p95 cho latency; FPS được đo
trên cùng khoảng thời gian. Power và thermal được thu trong bài chạy đủ dài để
quan sát trạng thái ổn định, không chỉ đo vài giây đầu.